%% -------------------------------------------------------
%% Kapitel 3: Systemarchitektur
%% -------------------------------------------------------
\chapter{Systemarchitektur}
\label{chap:architektur}

\section{Gesamtarchitektur}
\label{sec:gesamtarchitektur}

DoggoGO folgt einer klassischen \textbf{Client-Server-Architektur} mit drei Schichten:
einem Web-Frontend (Angular), einer nativen iOS-App und einem gemeinsamen Backend-Server,
der über eine REST-API kommuniziert. Die Datenhaltung erfolgt in einer PostgreSQL-Datenbank.
Abbildung~\ref{fig:gesamtarchitektur} illustriert die Systemkomponenten und deren Beziehungen.

\begin{figure}[H]
  \centering
  \begin{tikzpicture}[
      comp/.style={draw, rectangle, rounded corners=4pt,
                   minimum width=3.2cm, minimum height=0.9cm,
                   align=center, fill=gray!10, font=\small, text width=3cm},
      arr/.style={-{Stealth[length=5pt]}, thick},
      label/.style={font=\footnotesize\itshape, pos=0.65, above}
    ]
    \node[comp] (ios)     at (0, 1.2)  {iOS-App\\(Swift/SwiftUI)};
    \node[comp] (angular) at (0,-1.2)  {Angular\\Web-App};
    \node[comp] (api)     at (5, 0)    {Backend-API\\(ASP.NET Core 8)};
    \node[comp] (db)      at (10, 0)   {PostgreSQL\\Datenbank};

    \draw[arr] (ios)     -- node[label, yshift=5pt]{HTTP/JSON} (api);
    \draw[arr] (angular) -- node[label, below, yshift=-5pt]{HTTP/JSON} (api);
    \draw[arr] (api)     -- node[label, pos=0.4]{EF Core} (db);
  \end{tikzpicture}
  \caption{Gesamtarchitektur von DoggoGO}
  \label{fig:gesamtarchitektur}
\end{figure}

Die Entscheidung für eine gemeinsame Backend-API (anstelle von separaten Backends für Web
und iOS) hat folgende Vorteile:
\begin{itemize}
  \item Einmalige Implementierung der Geschäftslogik – keine Redundanz,
  \item einheitlicher Datenzugriff für beide Clients,
  \item einfache Erweiterbarkeit um weitere Clients (z.\,B.\ Android).
\end{itemize}

\section{Schichten des Backends}
\label{sec:schichten-backend}

Das Backend ist intern in drei Schichten gegliedert (Abbildung~\ref{fig:backend-schichten}):

\begin{description}
  \item[Controller-Schicht (Presentation Layer)]
    ASP.NET~Core~Controller nehmen HTTP-Anfragen entgegen, validieren Eingaben und
    delegieren die Verarbeitung an die Service-/Logikschicht.
  \item[Modell- und Datenbankschicht (Data Layer)]
    Entity~Framework~Core~9 übernimmt als ORM die Abbildung der C\#-Objekte auf
    PostgreSQL-Tabellen. \texttt{DoggoContext} ist der zentrale \texttt{DbContext}.
  \item[Hilfsdienste (Utility)]
    \texttt{BreedSeeder} importiert Hunderassen beim Start aus \texttt{breeds.json} in die
    Datenbank, sofern noch keine Einträge vorhanden sind.
\end{description}

\begin{figure}[H]
  \centering
  \begin{tikzpicture}[
      layer/.style={draw, rectangle, minimum width=5.5cm, minimum height=0.85cm,
                    align=center, font=\small, fill=gray!10, text width=5cm},
      util/.style={draw, rectangle, dashed, minimum width=3cm, minimum height=0.85cm,
                   align=center, font=\small, fill=white, text width=2.8cm},
      arr/.style={-{Stealth[length=5pt]}, thick}
    ]
    \node[layer] (ctrl) at (0, 3.0) {Controller-Schicht\\(HTTP-Requests, Routing)};
    \node[layer] (data) at (0, 1.5) {Datenbankschicht\\(DoggoContext / EF Core 9)};
    \node[layer] (pg)   at (0, 0)   {PostgreSQL};
    \node[util]  (seed) at (8, 1.5) {BreedSeeder\\(Utility)};

    \draw[arr] (ctrl) -- (data);
    \draw[arr] (data) -- (pg);
    \draw[arr, dashed] (seed) -- node[font=\footnotesize, above]{JSON-Seed} (data);
  \end{tikzpicture}
  \caption{Interne Schichten des Backends}
  \label{fig:backend-schichten}
\end{figure}

\section{Datenbankdesign}
\label{sec:datenbankdesign}

Das Datenbankschema besteht aus vier Entitäten (Abbildung~\ref{fig:er-diagramm}):

\begin{description}
  \item[\texttt{Owner}] Repräsentiert eine registrierte Nutzerin / einen registrierten
    Nutzer mit Anmeldedaten (E-Mail, Passwort-Hash, Passwort-Salt) und optionaler
    Telefonnummer.
  \item[\texttt{Dog}] Ein Hundeprofi mit Name und Alter, das einer Nutzerin / einem
    Nutzer (\texttt{OwnerId}) und einer Hunderasse (\texttt{BreedId}) zugeordnet ist.
  \item[\texttt{Breed}] Eine Hunderasse mit Bezeichnung und Beschreibung, verknüpft mit
    einem \texttt{LegalRestrictions}-Eintrag.
  \item[\texttt{LegalRestrictions}] Rechtskategorie (z.\,B.\ \texttt{spezielle\_hunderasse})
    und Liste konkreter Vorschriften (Leinen-/Maulkorbpflicht, ATP-Pflicht etc.) als
    PostgreSQL-Array.
\end{description}

\begin{figure}[H]
  \centering
  \begin{tikzpicture}[
      entity/.style={draw, rectangle, minimum width=2.8cm, minimum height=0.8cm,
                     align=center, font=\small\bfseries, fill=gray!15},
      rel/.style={draw, diamond, minimum width=1.4cm, minimum height=0.8cm,
                  align=center, font=\footnotesize, fill=white, inner sep=1pt},
      arr/.style={thick},
      card/.style={font=\footnotesize}
    ]
    % Entitäten
    \node[entity] (owner)  at (0, 0)    {Owner};
    \node[entity] (dog)    at (4, 0)    {Dog};
    \node[entity] (breed)  at (8, 0)    {Breed};
    \node[entity] (legal)  at (8,-2)    {LegalRestrictions};

    % Verbindungen mit Kardinalitäten
    \draw[arr] (owner) -- node[card, above]{1} node[card, below]{N} (dog);
    \draw[arr] (dog)   -- node[card, above]{N} node[card, below]{1} (breed);
    \draw[arr] (breed) -- node[card, right]{N} node[card, left]{1} (legal);

    % Annotationen
    \node[font=\footnotesize\itshape, below=0.1cm of owner] {Email (UNIQUE)};
    \node[font=\footnotesize\itshape, below=0.1cm of dog]   {CASCADE};
    \node[font=\footnotesize\itshape, right=0.1cm of breed] {RESTRICT};
  \end{tikzpicture}
  \caption{Entity-Relationship-Diagramm (DoggoGO-Datenbank)}
  \label{fig:er-diagramm}
\end{figure}

\subsection{Beziehungen und Löschverhalten}
\label{subsec:beziehungen}

Die Beziehungen zwischen den Entitäten wurden im \texttt{DoggoContext} mit \texttt{OnModelCreating}
explizit konfiguriert:

\begin{itemize}
  \item \textbf{Owner → Dog (1:N, Cascade Delete):} Wird eine Nutzerin / ein Nutzer gelöscht,
    werden alle zugehörigen Hundeprofile ebenfalls gelöscht. Dies entspricht dem erwarteten
    Verhalten beim Account-Delete.
  \item \textbf{Dog → Breed (N:1, Restrict Delete):} Eine Hunderasse kann nicht gelöscht
    werden, solange noch Hunde dieser Rasse existieren. Damit wird die referenzielle
    Integrität gewahrt.
  \item \textbf{Breed → LegalRestrictions (N:1, Restrict Delete):} Analog – ein
    Rechtsdatensatz kann nicht entfernt werden, solange Rassen darauf verweisen.
  \item \textbf{Owner.Email (Unique Index):} Verhindert mehrfache Registrierung mit
    derselben E-Mail-Adresse auf Datenbankebene.
\end{itemize}

\section{REST-API-Überblick}
\label{sec:api-ueberblick}

Die API folgt den REST-Prinzipien: Ressourcen werden über hierarchische URLs adressiert,
HTTP-Methoden spiegeln die CRUD-Operationen wider, und alle Antworten erfolgen im
JSON-Format. Tabelle~\ref{tab:api-endpunkte} listet die geplanten Endpunkte auf.

\begin{table}[H]
  \centering
  \caption{Übersicht der API-Endpunkte}
  \label{tab:api-endpunkte}
  \begin{tabularx}{\textwidth}{|l|l|Y|}
    \hline
    \textbf{Methode} & \textbf{Pfad} & \textbf{Beschreibung} \\
    \hline
    POST   & \texttt{/api/owners/register} & Neue Nutzerin / neuen Nutzer registrieren \\
    \hline
    POST   & \texttt{/api/owners/login}    & Anmelden, JWT erhalten \\
    \hline
    GET    & \texttt{/api/owners/\{id\}}   & Nutzerprofil abrufen (auth. erforderlich) \\
    \hline
    GET    & \texttt{/api/dogs}            & Alle eigenen Hunde abrufen \\
    \hline
    POST   & \texttt{/api/dogs}            & Neuen Hund anlegen \\
    \hline
    PUT    & \texttt{/api/dogs/\{id\}}     & Hundeprofi aktualisieren \\
    \hline
    DELETE & \texttt{/api/dogs/\{id\}}     & Hundeprofi löschen \\
    \hline
    GET    & \texttt{/api/breeds}          & Alle Hunderassen abrufen \\
    \hline
    GET    & \texttt{/api/breeds/\{id\}}   & Einzelne Rasse mit Rechtsvorschriften abrufen \\
    \hline
  \end{tabularx}
\end{table}

\section{Kommunikationsfluss}
\label{sec:kommunikationsfluss}

Abbildung~\ref{fig:sequenzdiagramm} zeigt den typischen Ablauf einer authentifizierten
Anfrage von der iOS-App an das Backend:

\begin{figure}[H]
  \centering
  \begin{tikzpicture}[
      lifeline/.style={draw, rectangle, minimum width=2.4cm, minimum height=0.7cm,
                       align=center, font=\small\bfseries, fill=gray!15, text width=2.2cm},
      msg/.style={-{Stealth[length=4pt]}, thin},
      ret/.style={-{Stealth[length=4pt]}, thin, dashed},
      lbl/.style={font=\footnotesize, above, midway}
    ]
    % Köpfe
    \node[lifeline] (ios) at (0,0)   {iOS-App};
    \node[lifeline] (api) at (6,0)   {Backend-API};
    \node[lifeline] (db)  at (12.5,0) {PostgreSQL};

    % Lebenslinien
    \draw[dotted] (ios.south) -- (0,-8.2);
    \draw[dotted] (api.south) -- (6,-8.2);
    \draw[dotted] (db.south)  -- (12.5,-8.2);

    % Nachrichten
    \draw[msg] (0,-1.1)    -- node[lbl]{POST /api/owners/login}       (6,-1.1);
    \draw[msg] (6,-1.85)   -- node[lbl]{SELECT Owner WHERE email}     (12.5,-1.85);
    \draw[ret] (12.5,-2.6) -- node[lbl]{Owner-Datensatz}              (6,-2.6);
    \draw[ret] (6,-3.35)   -- node[lbl]{200 OK + JWT}                 (0,-3.35);
    \node[font=\footnotesize\itshape, anchor=west] at (0.15,-4.0) {$\rightarrow$ Token in Keychain};
    \draw[msg] (0,-4.75)   -- node[lbl]{GET /api/dogs (Bearer Token)} (6,-4.75);
    \draw[msg] (6,-5.5)    -- node[lbl]{JWT validieren + SELECT Dogs} (12.5,-5.5);
    \draw[ret] (12.5,-6.25) -- node[lbl]{Dog-Datensätze}              (6,-6.25);
    \draw[ret] (6,-7.0)    -- node[lbl]{200 OK + Dogs\-Liste (JSON)}  (0,-7.0);
    \node[font=\footnotesize\itshape, anchor=west] at (0.15,-7.6) {$\rightarrow$ UI aktualisieren};
  \end{tikzpicture}
  \caption{Sequenzdiagramm – Authentifizierter API-Aufruf}
  \label{fig:sequenzdiagramm}
\end{figure}

\begin{enumerate}
  \item Die iOS-App sendet \texttt{POST /api/owners/login} mit Credentials im Request-Body.
  \item Das Backend prüft den HMAC-SHA512-Hash und gibt bei Erfolg ein signiertes JWT zurück.
  \item Die App speichert das Token sicher im iOS~\texttt{Keychain}.
  \item Für alle weiteren Anfragen wird das Token im \texttt{Authorization}-Header
        (\texttt{Bearer <token>}) mitgesendet.
  \item Der ASP.NET~Core Middleware-Stack validiert das Token vor jedem geschützten Endpunkt.
\end{enumerate}

\section{Sicherheitsarchitektur}
\label{sec:sicherheit}

Sicherheit ist ein Querschnittsthema, das alle Schichten des Systems betrifft. Für DoggoGO
wurden Maßnahmen auf drei Ebenen getroffen: Passwortspeicherung, Tokenbasierte
Authentifizierung und Autorisierung auf Datensatzebene.

\paragraph{Passwortspeicherung}
Passwörter werden niemals im Klartext gespeichert. Stattdessen wird für jede Nutzerin /
jeden Nutzer beim Registrieren ein individueller zufälliger Salt erzeugt
(\texttt{HMACSHA512.Key}), und das Passwort wird mit diesem Salt gehasht. Da jeder Salt
einzigartig ist, erzeugt dasselbe Passwort für unterschiedliche Nutzer\-innen und Nutzer
stets unterschiedliche Hashwerte. Damit sind vorberechnete Regenbogentabellen
(\emph{Rainbow Tables}) wirkungslos. Salt und Hash werden als \texttt{byte[]}-Felder in
der Datenbank persistiert.

\paragraph{JWT-Tokenfluss}
Nach erfolgreicher Anmeldung stellt der Server ein \textbf{JSON Web Token} (JWT) aus.
Das Token ist mit einem serverseitigen Secret (HMAC-SHA256) signiert und trägt als
Payload die \texttt{OwnerId} sowie die E-Mail-Adresse. Es hat eine konfigurierbare
Ablaufzeit (Standard: 8~Stunden). Der Client sendet das Token bei jeder nachfolgenden
Anfrage im \texttt{Authorization}-Header (\texttt{Bearer <token>}). Der ASP.NET~Core
Middleware-Stack validiert die Signatur und den Ablaufzeitstempel, bevor die
Controller-Methode aufgerufen wird – manipulierte oder abgelaufene Tokens werden mit
\texttt{401~Unauthorized} abgelehnt.

\paragraph{Autorisierung auf Datensatzebene}
Die \texttt{OwnerId} wird serverseitig aus den verifizierten JWT-Claims gelesen, nie aus
dem Request-Body. Dadurch ist es technisch unmöglich, dass eine angemeldete Nutzerin /
ein angemeldeter Nutzer Daten einer anderen Person liest oder verändert: Jede Datenbankabfrage
filtert implizit nach der eigenen \texttt{OwnerId}.

\paragraph{CORS}
Die API ist so konfiguriert, dass nur definierte Origins (Angular-Frontend und iOS-App)
Cross-Origin-Anfragen stellen dürfen. Damit werden unberechtigte Zugriffe von fremden
Webseiten auf die API unterbunden.

\section{Deployment-Architektur}
\label{sec:deployment}

Das gesamte Backend-Deployment basiert auf \textbf{Docker Compose}. Die
\texttt{compose.yaml}-Datei definiert zwei Services: \texttt{postgres} (offizielle
PostgreSQL-Image) und \texttt{backend} (das ASP.NET~Core-Anwendungs-Image). Beide Services
kommunizieren über ein gemeinsames Docker-Netzwerk; das Backend verwendet die Umgebungsvariablen
des Compose-Stacks für die Datenbankverbindung. Die PostgreSQL-Daten werden in einem benannten
Volume (\texttt{postgres\_data}) persistiert.
